% Source: https://artofproblemsolving.com/community/c3h2219004_some_identities
% Stable identity: aops-topic-2219004-q02
\textbf{Disproof.} For $n=6$: $\dbinom64=15$, but $1^2+2^2+3^2=14$. So the statement is false.

\textbf{All $n$ for which it holds.} Using $\displaystyle\sum_{i=1}^{k}i^2=\dfrac{k(k+1)(2k+1)}{6}$ with $k=n-3$,
\[\sum_{i=1}^{n-3}i^2=\dfrac{(n-3)(n-2)(2n-5)}{6},\qquad\dbinom n4=\dfrac{n(n-1)(n-2)(n-3)}{24}.\]
For $n\ge4$ the common factor $(n-2)(n-3)$ is positive, so the equation holds exactly when
\[\dfrac{n(n-1)}{24}=\dfrac{2n-5}{6}\iff n^2-n=8n-20\iff n^2-9n+20=0\iff(n-4)(n-5)=0.\]
So the equation holds only for $\boxed{n=4\text{ and }n=5}$ ($1=1$ and $5=1+4$). For $n\ge6$ we have $(n-4)(n-5)>0$, so $\dbinom n4$ is larger than the sum.
